\documentclass[12pt]{article}
\usepackage{amsmath,amssymb,amsfonts}
\usepackage[margin=1in]{geometry}

\begin{document}

\section*{Chapter 7, Problem 23}

\textbf{Problem:} Show that if the wave function $\psi$ of a particle satisfies Schr\"odinger's equation,
\[
H\psi = i\hbar\frac{\partial\psi}{\partial t},
\]
where $H$ is a self-adjoint operator, and if
\[
\int |\psi|^2\,d^3r\,\psi^* = 1 \quad \text{(1)}
\]
at $t = 0$, then at any subsequent time $t$, we will also have
\[
\int |\psi|^2\,d^3r = 1. \quad \text{(2)}
\]

Since in quantum mechanics $|\psi|^2$ is the probability density, $|\psi|^2\,d^3r$ gives the probability of finding the particle in the volume element $d^3r$, and Eqs.\ (1) and (2) are a statement of the conservation of total probability. If the total probability of finding the particle somewhere is initially equal to unity [Eq.\ (1)], it remains equal to unity at all later times. Clearly, the interpretation of $|\psi|^2$ as a probability density would not make sense if this were not the case. Note that such an interpretation would not hold for solutions to the heat (or diffusion) equation.

\bigskip
\textbf{Solution:}

We want to show that $\frac{d}{dt}\int|\psi|^2\,d^3r = 0$.

Compute:
\[
\frac{d}{dt}\int|\psi|^2\,d^3r = \int\Bigl(\frac{\partial\psi^*}{\partial t}\,\psi + \psi^*\frac{\partial\psi}{\partial t}\Bigr)\,d^3r.
\]

From Schr\"odinger's equation:
\[
\frac{\partial\psi}{\partial t} = \frac{1}{i\hbar}H\psi = -\frac{i}{\hbar}H\psi,
\]
and taking the complex conjugate (using $H = H^\dagger$, i.e., $H$ is self-adjoint):
\[
\frac{\partial\psi^*}{\partial t} = \frac{i}{\hbar}(H\psi)^* = \frac{i}{\hbar}H^\dagger\psi^* \quad\text{(formally)}.
\]

More precisely, $\frac{\partial\psi^*}{\partial t} = \frac{i}{\hbar}(H\psi)^*$.

Substituting:
\begin{align}
\frac{d}{dt}\int|\psi|^2\,d^3r &= \int\Bigl[\frac{i}{\hbar}(H\psi)^*\,\psi - \frac{i}{\hbar}\psi^*\,(H\psi)\Bigr]\,d^3r \\
&= \frac{i}{\hbar}\Bigl[\int(H\psi)^*\psi\,d^3r - \int\psi^*H\psi\,d^3r\Bigr] \\
&= \frac{i}{\hbar}\bigl[\langle H\psi|\psi\rangle - \langle\psi|H\psi\rangle\bigr].
\end{align}

Since $H$ is self-adjoint ($H = H^\dagger$):
\[
\langle H\psi|\psi\rangle = \langle\psi|H^\dagger\psi\rangle = \langle\psi|H\psi\rangle.
\]

Therefore:
\[
\frac{d}{dt}\int|\psi|^2\,d^3r = \frac{i}{\hbar}\bigl[\langle\psi|H\psi\rangle - \langle\psi|H\psi\rangle\bigr] = 0.
\]

Since the norm is conserved and initially equals 1, it remains 1 for all time.

\textbf{Comparison with the heat equation:} The heat equation $\partial T/\partial t = k\nabla^2 T$ has a real coefficient (no factor of $i$). Repeating the argument:
\[
\frac{d}{dt}\int T^2\,d^3r = 2k\int T\nabla^2 T\,d^3r = -2k\int|\nabla T|^2\,d^3r \le 0.
\]
So $\int T^2$ is \textit{not} conserved --- it decreases with time. This is why $T^2$ cannot be interpreted as a probability density. $\blacksquare$

\end{document}
